% Generated draft: requires mathematical and source review.
\documentclass[11pt]{article}
\usepackage[margin=25mm]{geometry}
\usepackage{fontspec}
\setmainfont{Latin Modern Roman}
\usepackage{amsmath,amssymb,mathtools,unicode-math}
\setmathfont{Latin Modern Math}
\usepackage{enumitem,needspace,xurl,hyperref}
\setlength{\parindent}{0pt}
\setlength{\parskip}{5pt}
\setlength{\emergencystretch}{3em}
\newcommand{\N}{\mathbb{N}}
\newcommand{\Z}{\mathbb{Z}}
\newcommand{\Q}{\mathbb{Q}}
\newcommand{\R}{\mathbb{R}}
\newcommand{\C}{\mathbb{C}}
\newcommand{\F}{\mathbb{F}}
\newcommand{\A}{\mathbb{A}}
\newcommand{\PP}{\mathbb{P}}
\newcommand{\E}{\mathbb{E}}
\newcommand{\eps}{\varepsilon}
\newcommand{\dd}{\,\mathrm{d}}
\begin{document}
\section*{UnsolvedMath ID: 30005644}
\textbf{Source problem code:} catalogue entry
\subsection{Definitions and mathematical statement}
Let \(p>0\) be a prime. For each \(i\in\{1,2\}\), let \(k_i\) be an algebraically closed field of characteristic \(p\), and let
\[
X_i^{\bullet}=(X_i,D_{X_i})
\]
be a pointed stable curve of the same type \((g,n)\) over \(k_i\). Here \(X_i\) is the underlying semistable curve of genus \(g\), and \(D_{X_i}\) is its set of marked points, with \(|D_{X_i}|=n\). In the usual stable-curve terminology, genus means arithmetic genus and marked points are distinct smooth points; stability requires every irreducible component of the normalization, of genus \(h\) and with \(r\) marked points or branches above nodes, to satisfy \(2h-2+r>0\).

Write \(\Pi_{X_i}^{\bullet}\) for the admissible fundamental group associated with \(X_i^{\bullet}\) in the supplied source. The excerpt does not give the construction or conventions for this group. For topological groups \(G,H\), the literal meaning of a continuous open homomorphism is a continuous group homomorphism \(\varphi:G\to H\) that sends open subsets of \(G\) to open subsets of \(H\). On this literal reading,
\[
\operatorname{Hom}_{\mathrm{op}}(G,H)
=\{\varphi:G\to H:\varphi\text{ is a continuous open group homomorphism}\}.
\]
The source calls the relevant maps open continuous homomorphisms, but the supplied excerpt does not specify basepoint conventions or whether maps are identified up to inner automorphisms. Consequently this displayed definition explains the literal reading without resolving those source conventions.

\emph{Source question.} Does there exist an anabelian interpretation of
\[
\operatorname{Hom}_{\mathrm{op}}
\bigl(\Pi_{X_1}^{\bullet},\Pi_{X_2}^{\bullet}\bigr)?
\]
The accompanying source discussion relates this question to deformation information of curves, moduli spaces of fundamental groups, and a homeomorphism conjecture. However, the supplied record does not define the desired geometric objects, equivalence relation, or correspondence constituting an anabelian interpretation. Thus the question remains a research target; the evidence does not determine a specific bijection, reconstruction theorem, or other precise conjectural assertion.
\subsection{Short English statement}
For two pointed stable curves of the same type over algebraically closed fields of a common positive characteristic, seek a geometric explanation of the open continuous homomorphisms between their admissible fundamental groups. The source connects this target with deformation information and moduli spaces, but leaves the precise correspondence and group conventions unspecified.
\subsection{Sources}
\begin{itemize}
\item Local UnsolvedMath registry, supplied record with permanent ID \(30005644\), source problem code \texttt{catalogue entry}, entitled ``Anabelian Interpretation of Homomorphisms between Fundamental Groups,'' in the Oberwolfach Reports Open Problems collection. No individual proposer is identified; the supplied proposed year is 2023.
\item Supplied report citation: \emph{MFO--RIMS Tandem Workshop: Arithmetic Homotopy and Galois Theory} (2024), Oberwolfach Reports, pp.~2377--2488, DOI \texttt{10.4171/owr/2023/42}. \url{https://doi.org/10.4171/owr/2023/42}. The record attributes the question to this report; the report contents have not been independently consulted for this draft. The cited year 2024 differs from the year 2023 in the DOI and supplied proposed-year metadata.
\item Supplied literature source: ``Yang, moduli spaces of fundamental groups.'' \url{https://arxiv.org/abs/2010.01806}. The supplied excerpt also refers to \texttt{[HYZ23]} in connection with deformation information, but supplies no complete bibliographic identification for that reference.
\item Supplied literature source: ``OWR report.'' \url{https://doi.org/10.4171/owr/2023/42}.
\item Literature assessment dated exactly \texttt{2026-08-22}: this is supplied dated curation, not independently verified. It labels the problem open and states that the moduli-space construction and several reconstruction results provide evidence for an anabelian explanation, while the extracted Hom-version question remains an open research problem. The background explicitly requests independent verification of time-sensitive status and cited claims.
\item The supplied registry labels the statement \texttt{corrected\_verified} and reports checking it against the official Oberwolfach report. These are supplied source-verification claims, not verification performed for this draft.
\end{itemize}
\end{document}
